\chapter{مقدمه و معرفی قالب}
\label{chap:intro}

به این قالب مدرن جزوه‌نویسی و کتاب‌نویسی با زی‌پرشین خوش آمدید. این قالب به صورت کاملاً حرفه‌ای و ماژولار طراحی شده است. هدف از طراحی این قالب، ارائه یک ساختار تمیز، قابل توسعه و زیبا برای نگارش متون علمی، ریاضی و برنامه‌نویسی است. 

\section{ساختار پروژه‌}
ساختار پوشه‌بندی این قالب به گونه‌ای طراحی شده تا مدیریت فایل‌ها به ساده‌ترین شکل ممکن انجام شود:
\begin{itemize}
    \item پوشه \E{config}: حاوی تنظیمات اصلی پروژه مانند بسته‌ها و فونت‌ها.
    \item پوشه \E{styles}: حاوی استایل‌های اختصاصی مانند محیط‌ها، جداول و کدها.
    \item پوشه \E{chapters}: حاوی فصل‌های مختلف کتاب یا جزوه.
    \item پوشه \E{codes}: حاوی سورس‌کدهای برنامه‌نویسی برای فراخوانی در متن.
    \item پوشه \E{figures}: حاوی فایل‌های اشکال و تصاویر.
    \item پوشه \E{tables}: حاوی کدهای مربوط به جداول.
\end{itemize}

در فصل‌های بعدی به صورت عملی با تمامی امکانات این قالب آشنا خواهید شد و نمونه‌هایی از خروجی‌های آن‌ها را مشاهده می‌کنید.
